\documentclass[11pt]{article}
\usepackage[margin=1in]{geometry}
\usepackage{amsmath,amssymb,booktabs}

\title{Depth Shortcut Predictor for DiT: Training Summary}
\date{}

\begin{document}
\maketitle

\section{Setup}
Let a DiT backbone have depth $L=12$, $P=256$ patch tokens, hidden size
$D=768$, and $T=50$ discrete training timesteps. For a noised latent
$x_t$ at timestep index $q\in\{0,\ldots,49\}$, a full DiT forward returns
\[
    \hat v,\quad Z_0,Z_1,\ldots,Z_L,\quad e_t^{\mathrm{DiT}},
\]
where $Z_0$ is the post-patch-embedding state and $Z_l\in\mathbb{R}^{B\times P\times D}$
is the state after block $l$.

All shortcut supervision is done in tokenwise direction space:
\[
    U_l=\operatorname{dir}(Z_l)
    =\frac{Z_l}{\|Z_l\|_2+\epsilon},
    \qquad U_l\in\mathbb{R}^{B\times P\times D}.
\]
The shortcut predictor receives a source direction, source layer, target
layer, and DiT timestep embedding:
\[
    Y_{a\rightarrow b}=P_\theta(U_a,a,b,t_q),\qquad
    \hat U_{a\rightarrow b}=\operatorname{dir}(Y_{a\rightarrow b}).
\]

\section{Conditioning}
For $0\leq a<b\leq L$, let $\Delta=b-a$. The predictor uses learned layer
embeddings of dimension $16$:
\[
    h_a=W_{\mathrm{state}}o_a,\quad
    h_b=W_{\mathrm{state}}o_b,\quad
    h_\Delta=W_\Delta o_\Delta .
\]
The depth condition is
\[
    h_{\mathrm{depth}}=[h_a,h_b,h_b-h_a,h_\Delta],
\]
and the predictor condition is produced by an MLP from
\[
    c_{\mathrm{in}}=
    [h_{\mathrm{depth}},\operatorname{sg}(e_t^{\mathrm{DiT}})].
\]
The stop-gradient prevents shortcut training from directly perturbing the
backbone timestep embedding through the condition path.

\section{Predictor Architecture}
The source direction is reshaped from $B\times256\times768$ to
$B\times16\times16\times768$. A ConvNeXt-style residual block is
\[
    h\leftarrow h+
    \mathrm{PW}_2\!\left(
    \mathrm{GELU}\!\left(
    \mathrm{PW}_1\!\left(
    \mathrm{DWConv}_{3\times3}(\mathrm{AdaLN}(h,c))
    \right)\right)\right),
\]
where
\[
    \mathrm{AdaLN}(h,c)=(1+\gamma(c))\mathrm{LN}(h)+\beta(c).
\]
The output projection predicts a residual direction update with a learnable
scalar output gate initialized at $\gamma_{\mathrm{out}}=0.05$:
\[
    Y_{a\rightarrow b}=U_a+\gamma_{\mathrm{out}}\Delta Y.
\]

\begin{center}
\begin{tabular}{lrrrrl}
\toprule
Version & Width & Blocks & Expansion & Params & Use case \\
\midrule
P-Tiny  & 256 & 2 & 2 & 1.3--1.8M & fast debug \\
P-Small & 384 & 4 & 2 & 3.5--4.2M & lightweight runs \\
P-Base  & 768 & 4 & 2 & 10--11M & main balanced model \\
P-Large & 768 & 6 & 2 & 15--17M & long skips / ceiling runs \\
\bottomrule
\end{tabular}
\end{center}

\section{Direct Shortcut Loss}
Sample a direct pair by first sampling a gap $d\sim\mathrm{Uniform}\{1,\ldots,L\}$
and then $a\sim\mathrm{Uniform}\{0,\ldots,L-d\}$, with $b=a+d$.
The direct loss is cosine distance to the stop-gradient target direction:
\[
    \mathcal{L}_{\mathrm{dir}}
    =
    1-
    \frac{1}{BP}\sum_{n=1}^{B}\sum_{i=1}^{P}
    \hat U_{a\rightarrow b}^{(n)}(i)^\top
    \operatorname{sg}(U_b^{(n)}(i)).
\]

\section{Bootstrap Depth Consistency}
Sample a triplet $0\leq a<b<c\leq L$ by sampling positive gaps
$d_1=b-a$ and $d_2=c-b$ uniformly and resampling if $d_1+d_2>L$.
Two short shortcuts are composed:
\[
    \hat U_{a\rightarrow b}
    =
    \operatorname{dir}(P_\theta(U_a,a,b,t_q)),
\]
\[
    \hat U_{a\rightarrow b\rightarrow c}
    =
    \operatorname{dir}(P_\theta(\hat U_{a\rightarrow b},b,c,t_q)).
\]
The target is the EMA teacher's long shortcut:
\[
    \hat U_{a\rightarrow c}^{\mathrm{ema}}
    =
    \operatorname{sg}\!\left[
    \operatorname{dir}(P_{\bar\theta}(U_a,a,c,t_q))
    \right].
\]
The bootstrap loss is
\[
    \mathcal{L}_{\mathrm{boot}}
    =
    1-
    \frac{1}{BP}\sum_{n,i}
    \hat U_{a\rightarrow b\rightarrow c}^{(n)}(i)^\top
    \hat U_{a\rightarrow c}^{\mathrm{ema},(n)}(i).
\]
The predictor teacher is updated after each optimizer step:
\[
    \bar\theta\leftarrow0.999\,\bar\theta+0.001\,\theta.
\]

\section{Total Objective}
The backbone is still trained by the original flow matching loss:
\[
    \mathcal{L}_{\mathrm{FM}}
    =
    \mathbb{E}\|\hat v-(x_0-x_1)\|_2^2 .
\]
The joint objective is
\[
    \boxed{
    \mathcal{L}
    =
    \mathcal{L}_{\mathrm{FM}}
    +0.5\,\mathcal{L}_{\mathrm{dir}}
    +0.25\,\mathcal{L}_{\mathrm{boot}}
    }.
\]
Targets $U_b$ and $P_{\bar\theta}$ are stop-gradient. Gradients from the
shortcut losses can flow through the source state $U_a$ into the DiT backbone.

\section{EMA Magnitude Calibration}
Magnitude is not used in the shortcut training losses. For skip inference
ablation only, maintain
\[
    R_{\mathrm{EMA}}\in\mathbb{R}^{(L+1)\times5\times P}.
\]
The timestep bin is $\tau(q)=\lfloor q/10\rfloor$. For each layer, bin, and
token:
\[
    R_{\mathrm{EMA}}[l,\tau,i]
    \leftarrow
    0.99\,R_{\mathrm{EMA}}[l,\tau,i]
    +0.01\,
    \mathbb{E}_{n:\tau(q_n)=\tau}\|Z_l^{(n)}(i)\|_2 .
\]
At skip inference from layer $a$ to $b$:
\[
    \tilde Z_b(i)=
    R_{\mathrm{EMA}}[b,\tau(q),i]\,
    \operatorname{dir}(P_\theta(\operatorname{dir}(Z_a),a,b,t_q))(i).
\]
The DiT then resumes from block $b+1$.

\section{Optimization and Metrics}
The DiT uses AdamW with backbone learning rate from the main config and
weight decay $0.01$, excluding biases, normalization parameters, and
embedding tables. The predictor uses AdamW with learning rate $10^{-4}$,
weight decay $0.1$, betas $(0.9,0.999)$, and gradient clipping at $1.0$.

Log the training losses and cosine similarities:
\[
    \cos_{\mathrm{dir}}=
    \frac{1}{BP}\sum_{n,i}
    \hat U_{a\rightarrow b}^{(n)}(i)^\top U_b^{(n)}(i),
\]
\[
    \cos_{\mathrm{boot}}=
    \frac{1}{BP}\sum_{n,i}
    \hat U_{a\rightarrow b\rightarrow c}^{(n)}(i)^\top
    \hat U_{a\rightarrow c}^{\mathrm{ema},(n)}(i).
\]
For inference ablations, also report restored-state relative error, output
deviation, FID, wall-clock speedup, and the FID--speed curve.

\section{Skip-FID Evaluation}
In addition to full-backbone FID, evaluate a skip sampler that repeatedly
jumps by two state indices:
\[
    1\rightarrow3,\quad 3\rightarrow5,\quad 5\rightarrow7,\quad \ldots
\]
At each source layer $a$, the predictor estimates only direction:
\[
    \hat U_{a\rightarrow a+2}
    =
    \operatorname{dir}(P_{\bar\theta}(\operatorname{dir}(Z_a),a,a+2,t_q)).
\]
Before continuing the DiT, restore the target-layer magnitude with the L2 EMA:
\[
    \tilde Z_{a+2}(i)=
    R_{\mathrm{EMA}}[a+2,\tau(q),i]\,
    \hat U_{a\rightarrow a+2}(i).
\]
The sampler then skips the true intervening blocks and resumes from the block
after the restored target state. Report these metrics separately, e.g.
$\mathrm{FID}_{\mathrm{skip}}$ and $\mathrm{sFID}_{\mathrm{skip}}$, so the
quality--speed tradeoff is visible next to full-backbone FID.

\section{Experimental Plan}
Run the method in four stages so implementation bugs, capacity limits, and
skip-quality tradeoffs are separated.

\paragraph{Stage 1: Pipeline debug.}
Train with P-Tiny for short runs. The goal is only to verify that the joint
loss is numerically stable, EMA teacher updates, $R_{\mathrm{EMA}}$ fills all
five timestep bins, and the raw predictor norm remains reasonable:
\[
    \mathbb{E}_{n,i}\|Y_{a\rightarrow b}^{(n)}(i)\|_2,\qquad
    \min_{n,i}\|Y_{a\rightarrow b}^{(n)}(i)\|_2,\qquad
    \max_{n,i}\|Y_{a\rightarrow b}^{(n)}(i)\|_2 .
\]
Since $\|U_a(i)\|_2=1$, these values should not collapse near zero or grow to
large magnitudes before direction normalization.

\paragraph{Stage 2: Main model.}
Train P-Base as the balanced setting for the primary result. Track
$\mathcal{L}_{\mathrm{FM}}$, $\mathcal{L}_{\mathrm{dir}}$,
$\mathcal{L}_{\mathrm{boot}}$, $\cos_{\mathrm{dir}}$, and
$\cos_{\mathrm{boot}}$. Evaluate both full-backbone FID and skip-FID at the
same checkpoint cadence.

\paragraph{Stage 3: Capacity sweep.}
If P-Base underfits long gaps, run P-Large with the same optimizer settings.
Compare long-gap direct cosine and skip-FID against P-Base for skips such as
$1\rightarrow3$, $3\rightarrow5$, and longer diagnostic pairs like
$0\rightarrow6$ or $0\rightarrow12$.

\paragraph{Stage 4: Quality--speed curve.}
For selected checkpoints, report full FID/sFID and skip FID/sFID together
with wall-clock sampling time. The useful operating point is the largest
speedup whose FID degradation remains acceptable.

\end{document}
